% !tp-example Basic usage
\documentclass{article}
\begin{document}
\wcDiv[data-tp-reference-scope]{
  \wcTpListof[selector=figure]{}
  \wcFigure{
    A diagram of the water cycle.

    \wcFigcaption{The water cycle}
  }
  \wcFigure{
    A diagram of a food chain.

    \wcFigcaption{A food chain}
  }
}
\end{document}

% !tp-example Javascript API
\documentclass{article}
\begin{document}
\wcScript[type=module]|
import "/components/listof/listof.js";

const example = document.createDocumentFragment();

const element1 = document.createElement("div");
element1.setAttribute("data-tp-reference-scope", "");
element1.append(document.createTextNode("\n  "));

const element2 = document.createElement("tp-listof");
element2.setAttribute("selector", "figure");
element1.append(element2);
element1.append(document.createTextNode("\n  "));

const element3 = document.createElement("figure");

const element4 = document.createElement("p");
element4.append(document.createTextNode("A diagram of the water cycle."));
element3.append(element4);

const element5 = document.createElement("figcaption");
element5.append(document.createTextNode("The water cycle"));
element3.append(element5);
element1.append(element3);
element1.append(document.createTextNode("\n  "));

const element6 = document.createElement("figure");

const element7 = document.createElement("p");
element7.append(document.createTextNode("A diagram of a food chain."));
element6.append(element7);

const element8 = document.createElement("figcaption");
element8.append(document.createTextNode("A food chain"));
element6.append(element8);
element1.append(element6);
element1.append(document.createTextNode("\n"));
example.append(element1);

document.body.append(example);
|
\end{document}

% !tp-example Notes, bibliography and glossary
\documentclass{article}
\begin{document}
\wcDiv[data-tp-reference-scope]{
  Read a note \rawhtml*|<tp-ref href="^method"></tp-ref>|, consult a source \wcTpRef*[href=@book]{}, or look up \rawhtml*|<tp-ref href="%ecosystem"></tp-ref>|. Compare another note \rawhtml*|<tp-ref href="^weather"></tp-ref>|, then read the first note again \rawhtml*|<tp-ref href="^method"></tp-ref>|.

  \wcTpNote[ref=method,title=Method]{
    Observe the habitat at \textbf{three different times} of day.

    \begin{itemize}
      \item Record the weather.
      \item Compare your observations.
    \end{itemize}
  }
  \wcTpBiblio[ref=book,title=Book]{
    Alex Example. \emph{Field observation handbook}. Example Press, 2026.

  }
  \wcTpGlossary[ref=ecosystem,title=Ecosystem]{
    A community of organisms interacting with their physical environment.

  }
  \wcTpNote[ref=weather]{
    Record the temperature and cloud cover.

  }
  \wcTpBiblio[ref=atlas]{
    Alex Example. \emph{A field atlas}. Example Press, 2025.

  }
  \wcTpGlossary[ref=community]{
    A group of interacting populations.

  }
  \subsection*{Notes}
  \wcTpListof[selector=tp-note]{}
  \subsection*{Bibliography}
  \wcTpListof[selector=tp-biblio]{}
  \subsection*{Glossary}
  \wcTpListof[selector=tp-glossary]{}
}
\end{document}
